\metatitle{Koornwinder polynomials}
\metadescription{Symmetric, nonsymmetric, and relative Koornwinder polynomials,
their double-affine Hecke algebra definition, specializations, and
combinatorial formulas.}
\metakeywords{Koornwinder polynomial, Macdonald-Koornwinder polynomial,
Askey-Wilson polynomial, double affine Hecke algebra, signed permutation,
alcove walk, set-valued tableau}


\section[koornwinderPolynomial]{Koornwinder polynomials}

\begin{polydata}{koornwinderPolynomial}
  Name     & Koornwinder polynomials \\
  Space    & Other \\
  Basis    & True \\
  Rating   & 3 \\
  Bib      & Koornwinder1992BC \\
  Year     & 1992 \\
  Symbol   & \(P_\lambda^{\mathrm K}(\xvec;q,t;a,b,c,d)\) \\
  Category & Macdonald \\
\end{polydata}

The \defin{Koornwinder polynomials}\label{koornwinderPolynomialDefinition},
also called \defin{Macdonald--Koornwinder polynomials}, are the Macdonald
polynomials for the nonreduced affine root system
\((C_n^\vee,C_n)\) \cite{Koornwinder1992BC}.  They form the universal
classical-type Macdonald family:
the Macdonald polynomials for every classical affine root system arise by
specializing their parameters \cite{ColmenarejoGagnonRam2026x}.
For \(n=1\), they are the Askey--Wilson polynomials.

Unlike type \(A\) \hyperref[macdonaldP]{Macdonald \(P\) polynomials},
Koornwinder polynomials are Laurent polynomials.  They depend on six
parameters, customarily written \(q,t,a,b,c,d\).  Other common conventions use
square-root parameters
\(q^{1/2},t^{1/2},t_0^{1/2},u_0^{1/2},t_n^{1/2},u_n^{1/2}\).


\subsection[koornwinderDefinition]{Eigenfunction and triangularity definition}

Let

\[
  W_{\mathrm{fin}}=\symS_n\ltimes(\setZ/2\setZ)^n
\]

be the signed permutation group.  It permutes the variables and replaces
individual \(x_i\) by \(x_i^{-1}\).  For a partition
\(\lambda=(\lambda_1\geq\dotsb\geq\lambda_n\geq0)\), let

\[
  m_\lambda^{BC}(\xvec)
  \coloneqq
  \sum_{\alpha\in W_{\mathrm{fin}}\lambda}\xvec^\alpha
\]

be the corresponding orbit sum, with each distinct orbit element included
once.

The symmetric Koornwinder polynomial \(P_\lambda^{\mathrm K}\) is the unique
\(W_{\mathrm{fin}}\)-invariant Laurent polynomial which is triangular,

\[
  P_\lambda^{\mathrm K}
  =
  m_\lambda^{BC}
  +\sum_{\mu\lhd\lambda}c_{\lambda\mu}(q,t;a,b,c,d)m_\mu^{BC},
\]

and is a joint eigenfunction of the commuting Koornwinder difference
operators.  Equivalently, the family is triangular and orthogonal for the
Koornwinder inner product.  This is the type \(BC\) analogue of the
\hyperref[macdonaldPEigenvectors]{eigenfunction characterization of Macdonald
polynomials}.

\begin{example}
The first orbit sums make the shape of the family visible:
\[
  P_{\varnothing}^{\mathrm K}=1,
  \qquad
  P_{(1)}^{\mathrm K}
  =\sum_{i=1}^n(x_i+x_i^{-1})+c,
\]
where \(c\) is a rational function of the six parameters.  Similarly, the
leading orbit of \(P_{(1,1)}^{\mathrm K}\) is
\[
  \sum_{1\leq i<j\leq n}
  (x_ix_j+x_ix_j^{-1}+x_i^{-1}x_j+x_i^{-1}x_j^{-1}).
\]
Lower orbit sums are determined by the eigenvalue equations.
\end{example}


\subsection[koornwinderNonsymmetric]{Electronic and relative polynomials}

The nonsymmetric version is indexed by \(\mu\in\setZ^n\).  Following the
terminology of \name{Laura Colmenarejo} and \name{Arun Ram}, it is useful to
call \(E_\mu^{\mathrm K}\) the \defin{electronic Koornwinder polynomial} and
\(P_\lambda^{\mathrm K}\) the \defin{bosonic Koornwinder polynomial}; compare
the corresponding terminology on the
\hyperref[macdonaldP]{Macdonald page} \cite{ColmenarejoRam2022x}.

The electronic polynomial is characterized by
\[
  Y_jE_\mu^{\mathrm K}=\gamma_{\mu,j}E_\mu^{\mathrm K},
  \qquad
  [\xvec^\mu]E_\mu^{\mathrm K}=1,
\]
where \(Y_1,\dotsc,Y_n\) are the commuting Cherednik operators for the
double-affine Hecke algebra of type \((C_n^\vee,C_n)\), and the eigenvalues
\(\gamma_{\mu,j}\) are explicit monomials in the six parameters.

For a signed permutation \(z\in W_{\mathrm{fin}}\), the
\defin{relative Koornwinder polynomial} \(E_\mu^z\) is a normalized Hecke
operator translate of \(E_\mu^{\mathrm K}\), with leading monomial
\(\xvec^{z\mu}\).  These are the type \(BC\) analogues of
\hyperref[macdonaldEperm]{permuted-basement Macdonald polynomials};
they also occur as open-boundary ASEP polynomials.  For dominant
\(\lambda\), symmetrization gives
\[
  P_\lambda^{\mathrm K}
  =\sum_{w\in W_{\mathrm{fin}}}E_\lambda^w.
\]


\subsection[koornwinderCombinatorialFormulas]{Combinatorial formulas}

\name{Laura Colmenarejo}, \name{Lucas Gagnon}, and \name{Arun Ram} give three
parallel formulas for all relative Koornwinder polynomials
\cite{ColmenarejoGagnonRam2026x}:

\begin{itemize}
\item a creation formula written with the
  \hyperref[dividedDifferenceOperators]{divided-difference operators} familiar
  from \hyperref[schubert]{Schubert calculus};
\item an alcove-walk formula reparametrized by uncompressed set-valued
  tableaux; and
\item a compressed set-valued tableau formula obtained by
  \enquote{across-the-\(0\)-gap} and \enquote{around-the-end} compression.
\end{itemize}

The relative family is essential here: it plays the same role as a variable
basement in type \(A\).  The box-greedy reduced word and its coroot sequence
replace the arm-and-leg data that organize the usual Macdonald tableau
formula.

\name{Sylvie Corteel}, \name{Olya Mandelshtam}, and
\name{Lauren Williams} give a different formula using rhombic staircase
tableaux \cite{CorteelMandelshtamWilliams2024}.  The precise combinatorial
relation between the compressed set-valued tableaux and the rhombic
staircase tableaux is currently open.


\subsection[koornwinderSpecializations]{Specializations and nearby families}

\begin{itemize}
\item For \(n=1\), \(P_{(r)}^{\mathrm K}\) is an Askey--Wilson polynomial.
\item Classical-type Macdonald polynomials are obtained by parameter
  specialization.
\item Hall--Littlewood limits of types \(B\) and \(C\) arise from the
  corresponding \(q=0\) specializations.
\item Relative Koornwinder polynomials specialize to families related to the
  open-boundary asymmetric exclusion process.
\end{itemize}

See the pages on \hyperref[rootSystem]{root systems},
\hyperref[coxeterGroups]{Coxeter groups}, and
\hyperref[macdonaldE]{nonsymmetric Macdonald polynomials} for the surrounding
type-\(A\) and Weyl-group background.
